\documentclass[12pt]{article}

\begin {document}
\section{NEURODEGENERACIÓN: Alzheimer,\\Parkinson y ELA}
	
	Este libro, plantea el objetivo de transmitir de manera sencilla y clara conceptos básicos de la neurodegeneración y de las principales enfermedades neurdegenerativas, intentado simplificar algo tan complejo como el cerebro humano. 
	
	\subsection{Definición y concepto}
	
	Inicialmente, se establece la definición básica de envejecimiento como principal factor de riesgo para el proceso degenerativo. A continuación, se trata superficialmente los mecanismos importantes tanto en el envejecimiento como en la degeneración, y se expone la definición del mismo en lo relacionado con el cerebro.\\
	
	La definición del envejecimiento se centra en la pérdida de vitalidad y el aumento progresivo de la vulnerabilidad y la fragilidad.  Este proceso biológico universal, progresivo, determinado-regulado, supresor heterogéneo e irreversible, conlleva una progresiva e irreversible pérdida de la integridad fisiológica.\\
	
	En la sociedad actual, la esperanza de vida tiende a un incremento progresivo. En los países desarrollados está en aumento, tanto la Edad Media de la población como el porcentaje de personas dentro del rango entre 85 y 100 años. Esto pone en el punto de mira la búsqueda del retraso de la aparición de enfermedades asociadas al envejecimiento.\\
	
	Gracias a varias investigaciones, se ha descubierto que las dietas calóricas son un importante factor de disminución de las expectativas de vida. Mediante mecanismos de reprogramación que restringen la cantidad de calorías en la alimentación, se consiguió incrementar la longevidad en numerosas especies y mejorar diversas actividades metabólicas. \\

	Por otra parte, ganó importancia la identificación de genes encargados de modular la longevidad. En algunos modelos animales se identificaron elementos modificadores asociados al aumento de las expectativas de vida, como por ejemplo:
	\begin{itemize}
		\item Disminución del factor de crecimiento insulina-like (IGF.1)
		\item Disminución de la proteína mTOR
		\item Aumento de la actividad de las siete sirtuinas conocidas, sobre todo de la sirtuina 1 (SIRT1) mediante la ingesta de resveratrol
	\end{itemize}
	
	Otro factor fundamental en esta línea de investigación es la epigenética. Actualmente se la define como el estudio de cambios heredables en la función genética que se producen sin ningún cambio en la secuencia del ADN. Es una ciencia que estudia la forma, como los factores ambientales y de estilo de vida que influyen en la expresión final de ciertos genes. Por esto mismo, el 
	epigenoma cambia durante toda la vida en función de estas variables, siendo muy flexible y conllevando grandes implicaciones en el campo de la salud.\\
	
	Muchas de las enfermedades debidas al envejecimiento, como algunos tipos de demencia, presentan un componente epigenético. Al envejecer, se pierden grupos metilos, las neuronas disminuyen o dejan de producir determinados neurotransmistores.\\
	
	La posibilidad de lograr un rejuvenecimiento con los cambios epigenéticos, lleva a hablar de la reprogramación celular. Esto se el método para la modificación del epigenoma.	Mediante la combinación de cuatro genes, llamados los factores de Yamanaka pudo conseguirse la detención o reversión del envejecimiento. Esto fue descubierto por Shanya Yamanaka en el 2006, al convertir una célula adulta en una célula madre pluripotente. Para evitar que haya más probabilidad de desarrollar células tumorales, se llegó a la conclusión de no permitir la finalización del proceso de transformación.\\
	
	Los factores ambientales modulan mecanismos epigenéticos que a su vez, modulan la capacidad funcional de las células madre.  Un decline de la funcionalidad de estas, lleva a la perdida de homeostasis. Recientemente se demostró la posibilidad de obtener células funcionalmente completas a partir de células senescentes. Esto era algo inimaginable hace no tanto en la reprogramación celular, por lo que esto muestra una evidencia adicional a la reversión del envejecimiento celular. Esto podría concluir en mejoras terapéuticas para personas con progeria.\\
	
	Por otro lado, Se sugiere la existencia de un reloj metabólico que controla el envejecimiento. Este Se basa en 9 sellos distintivos del envejecimiento, los cuales los podemos clasificar en 3 categorías. En La primera, nos encontramos con las principales causas del daño molecular, Como la inestabilidad genómica Basada en la rotura del equilibrio entre los daños y la capacidad de reparación; Alteraciones epigenéticas que cambian la estructura del ADN; pérdida de proteostasis Y desgaste del telómero. La segunda categoría están los componentes beneficiosos a niveles bajos y dañinos a niveles altos. Entre ellos se encuentran la disfunción mitocondrial, la senescencia celular, o la detección desregulada de nutrientes. Por último, Se encuentran los elementos con los que desaparece la capacidad de autorregulación, Como el cansancio de células madre y las alteraciones en la comunicación intercelular.\\

	La neurodegeneración consiste en el deterioro de células y tejidos, haciendo que estos pierdan de forma gradual la capacidad para realizar sus funciones. Por otra parte, esta la regeneración, proceso donde se restablecen células, tejidos y órganos. En cambio, las neuronas son células postmitóticas no tienen la habilidad de dividirse ni remplazarse, por lo que envejecen por la acumulación de daños con el paso del tiempo. Esto lleva a una pérdida de neuronas que afecta sobre todo al hipocampo y a los cingulados anterior y posterior, zonas especialmente afectadas en las enfermedades neurodegenerativas. También supone una pérdida de dendritas y sinapsis. 
	Hay varios mecanismos encargados de inducir la neurodegeneración. Entre ellos encontramos factores diversos como pueden ser genéticos y ambientales, las dinámicas anormales de las proteínas o el daño oxidativo.   Todo esto junto con otros sistemas llevan a la muerte celular.\\
	
	Intentando aclarar el mecanismo de degeneración de las células, algunos científicos llegaron a la conclusión de que la agregación de proteínas alteradas o deficientemente plegadas eran una de las principales causas de la muerte neuronal. Las distintas enfermedades neurodegenerativas tienen esta característica en común. Mediante distintos estudios, se consiguieron identificar ciertas alteraciones de proteínas correspondientes a diferentes enfermedades. Como, por ejemplo, el Alzheimer se asocia a la proteína B amiloide y el Parkinson a los depósitos de alfa-sinucleína.  Aún así, siguen abiertas ciertas preguntas, como el por qué de que algunas proteínas se plieguen de forma anormal y por qué algunas son más propensas a la agregación que otras.
	Otro concepto muy importante y bastante reciente es la resiliencia. Esta consiste en la capacidad de los individuos a ajustarse a situaciones adversas o, en el ámbito de la física, es la capacidad de ciertos materiales para recuperarse de ciertas cargas. Los factores determinantes de la resiliencia pueden ser determinantes para las investigaciones de la neurodegeneración.

	\subsection{Bases moleculares y celulares de la degeneración}
	
	El conocimiento de las bases moleculares y celulares de la neurodegeneración son esenciales para desarrollar nuevos métodos diagnósticos así como nuevas terapias. En el libro se aborda la compleja interacción entre los diferentes procesos biológicos y fisiológicos que pueden influir en la patogénesis de diferentes enfermedades. Se ha hecho un enorme esfuerzo en estudiar los mecanismos que regulan la muerte celular, ya que esto permitirá avanzar en nuevas intervenciones terapéuticas que detengan o enlentezcan el proceso neurodegenerativo.\\
	
	Se describen tres tipos de muerte celular: la apoptosis o “muerte celular programada”, la necrosis que se produce como consecuencia de la acción de toxinas, generando una caída de ATP y la autofagia. Este último, es un mecanismo empleado por las células eucariotas para degradar sus componentes con vistas a la renovación y recambio de estos. Se concluye que en la neurodegeneración la autofagia está descontrolada. Después de analizar este mecanismo en las principales enfermedades degenerativas, se concluye que se encuentran involucrados de una manera directa defectos en su funcionamiento.\\
	
	Otro de los procesos objeto de estudio es la oxidación y el estrés oxidativo. Esto es debido a nuestra naturaleza como seres vivos que consumen oxígeno y que están expuestos de manera permanente a subproductos tóxicos conocidos genéricamente como  “radicales libres de oxígeno” y “radicales libres de hidrógeno”. Uno de los radicales libre de mayor interés en medicina es el óxido nítrico, se estima que tiene un papel protagonista indirecto en la neurodegeneración a través de un derivado denominado “peroxinitrito”. Los resultados de diferentes estudios demuestran que la generación de prooxidantes es un factor que influye en la velocidad de envejecimiento, el cual está relacionado directamente con el daño mitocondrial.\\
	
	Por tanto, en el origen de las enfermedades neurodegenerativas se encuentran involucrada la elevada producción de especies reactivas de oxígeno (ROS) y especies reactivas de nitrógeno (RNS) asociadas a un mal funcionamiento de la mitocondria neuronal. El daño en el ADN, como consecuencia del estrés oxidativo, genera desarrollo de mutaciones en el mismo. En definitiva, se genera una relación multidireccional entre estrés oxidativo, mal funcionamiento mitocondrial, incremento de apoptosis y disminución de la autofagia, cuyo resultado es el desarrollo de la neurodegeneración.\\
	
	Un último elemento a analizar es la neuroinflamación como respuesta inflamatoria del sistema nervioso central a cualquier tipo de agresión o lesión sobre este. En este punto es relevante las funciones que tiene la glía, células que desempeñan un papel esencial de soporte a las neuronas y a ala fisiología del sistema nervioso. Aunque el proceso inflamatorio es un mecanismo de defensa del organismo contra la infección, la pérdida de su regulación hace que contribuya a una condición patológica, pues mantiene un proceso de destrucción celular progresivo. Esto nos lleva a que las respuestas inflamatorias se mantienen activadas durante el estado de estrés oxidativo, lo que lleva a las neuronas a un círculo vicioso de muerte y destrucción.	
	
	\subsection{Enfermedades neurodegenerativas}
	
	Las enfermedades degenerativas tienen cada vez mayor protagonismo en nuestra sociedad debido al envejecimiento de la población. Se definen como aquellas patologías heredadas o adquiridas, en las que se produce una disfunción progresiva del sistema nervioso central. Existen más de 600 enfermedades neurodegenerativas, pero destacan la enfermedad de Alzheimer (EA), la enfermedad de Parkinson (EP), la enfermedad de Huntington (EH) y la esclerosis lateral amiotrófica (ELA).
	
		\subsubsection{Enfermedad de Alzheimer}
		
		La EA es actualmente un problema de primer orden y una prioridad para los sistemas de salud. Hay varios estudios sobre su prevalencia, el estudio Delphi prevé que en el 2040 habrá más de 80 millones de personas afectadas por la enfermedad, existiendo unanimidad en que la prevalencia se dobla cada cinco años. Siendo una de las principales causas de muerte no se puede prevenir, curar ni siquiera ralentizar su avance.
		La edad es el principal factor de riesgo para su desarrollo, y presenta mayor incidencia en el género femenino en edades avanzadas, también un abaja escolarización aumenta el riesgo. Existen multitud de estudios sobre los diferentes factores de riesgo, una relación compleja difícil de interpretar debido a la alta variabilidad. La mayoría de los casos se considera que no tienen origen genético. Sin embargo, existen pruebas de factores genéticos que puedan conllevar una mayor predisposición a padecer la enfermedad, confirmándose el genotipo apolipoproteçina E (ApoE) como factor de riesgo en múltiples estudios.
		En la corteza cerebral, la sustancia blanca contiene los axones (responsables de las conexiones entre neuronas). La pérdida de sinapsis podría constituir la vía final común de la patología, ya que su número disminuye con la evolución de la enfermedad.\\
		
		Existe la necesidad de definir el estado entre normalidad y demencia y para ello han surgido en los últimos años diferentes criterios, cobrando especial relevancia en el mundo científico el término “deterioro cognitivo leve”. Los pacientes se quejan básicamente de pérdida de memoria, pero pueden manifestar otras dificultades. Las técnicas de neuroimagen son de indudable interés clínico, de forma complementaria a las manifestaciones clínicas. El objetivo es identificar a aquellos pacientes con mayor riesgo de desarrollar una demencia, e instalar un acercamiento terapéutico de manera temprana.
		La demencia es un proceso progresivo. Al déficit de memoria se van añadiendo dificultades del lenguaje, la desorientación y la desorganización de las actividades cotidianas. Entre otras manifestaciones clínicas se encuentra la apraxia y la agnosia y la afasia en sus diferentes manifestaciones. Los síntomas psicológicos y conductuales afectan a los pacientes con demencia, con una prevalencia del 90\%, entre los que cabe destacar la apatía, la agresividad, así como las alucinaciones y la alteración del ciclo vigilia-sueño. Es relevante ubicar al paciente en el estadio evolutivo para programar adecuadamente su tratamiento.\\
		
		El diagnóstico sigue siendo fundamentalmente clínico. Aunque en la ayuda al diagnóstico se empiezan a incorporar diversos biomarcadores, así como técnicas de imagen como la resonancia nuclear magnética, la tomografía por emisión de fotón único (SPECT) o la tomografía por emisión de positrones con fluordeoxiglucosa (PET-FDG). Esta última ha demostrado su utilidad para distinguir pacientes con EA de otras demencias. Otras técnicas de imagen pueden ser útiles para el diagnóstico precoz y el seguimiento de la enfermedad.\\
		
		La intervención terapéutica debe contemplar tratamiento farmacéutico, intervención no farmacológica,  y la actuación sobre familiares/cuidadores. El tratamiento farmacológico se centra en paliar la degeneración neuronal, potenciar la actividad de las neuronas no afectadas y aliviar los síntomas de la enfermedad. El tratamiento se centra en dos grupos de fármacos: los inhibidores de acetilcolinesterasa (IACE) y los antagonistas de los receptores N-Metil-D-Aspartano tipo memantina. Una línea abierta de investigación es la búsqueda de tratamientos que sean modificadores de la historia natural de la enfermedad, destacando el dirigido contra la proteína beta amiloide o la inmunoterapia activa. En lo que respecta a las intervenciones no farmacológicas, cabe destacar: programas de orientación a la realidad, terapia de reminiscencia, terapia de validación, estimulación cognitiva y estimulación magnética. Entre las medidas aplicadas al cuidador, señalar las orientadas a los programas de información y formación, los equipos de soporte domiciliario y las unidades de apoyo sociosanitario.\\
		
		En términos de coste económico, se considera que las demencias son la tercera causa más costosa, tras el cáncer y las enfermedades cardiovasculares. En más del 80\% de los casos, el cuidado lo realiza la familia. El diagnóstico precoz y el empleo de alternativas no farmacológicas podrían reducir el coste total de las demencias.		
		
		\subsubsection{Enfermedad de Parkinson (EP)}
		
		Es un trastorno neurodegenerativo crónico de curso progresivo y evolución prolongada, caracterizado por la pérdida de neuronas dopaminérgicas. Se trata de la segunda enfermedad neurodegenerativa más frecuente en mayores de 65 años, tras la enfermedad del Alzheimer y afecta entre el uno y el dos por ciento de la población. En la mayoría de los casos es de tipo esporádico y de causa desconocida.\\
		
		Las manifestaciones iniciales suelen pasar desapercibidas. Un síntoma común suele ser el temblor en reposo y de tipo asimétrico que afecta a la mano, aunque pueden aparecer otros síntomas iniciales. Existen una serie de síntomas no motores cuya presencia puede poner sobre aviso: trastornos en el olfato, alteraciones en la fase REM de sueño y el estreñimiento. El diagnóstico precoz a través de biomarcadores es primordial, desde técnicas de neuroimagen hasta marcadores genéticos, moleculares y tisulares. Las principales manifestaciones clínicas son: Bradicinesia. Temblor, Rigidez e inestabilidad postural. A estos se les pueden unir otros síntomas motores de manera variable: Disfagia, Síndrome de piernas inquietas y alteraciones del lenguaje. Entra las manifestaciones no motoras se encuentran: depresión, ansiedad, apatía,  en algunos casos problemas de vejiga o de estreñimiento, excesiva producción de saliva o inadecuada deglución, así como problemas en la piel. Un diagnóstico temprano puede ser difícil por falta de biomarcadores con suficiente nivel de sensibilidad y especificidad. El infradiagnóstico se sitúa en un 24 \% y es frecuente un diagnóstico erróneo. Pueden ser de ayuda pruebas de neuroimagen funcional.\\
		
		No hay cura para la EP, pero distintos medicamentos pueden proporcionar un significativo alivio en fases iniciales y debe continuarse durante casi toda la vida del paciente. La levopoda es el tratamiento sintomático más eficaz, permitiendo a los pacientes alargar el período de tiempo en el que poder llevar una vida relativamente normal. Otros fármacos, aunque menos eficaces en controlar la rigidez y la bradicinesia son la bromocriptina y la pergolida. La Selegilina se está utilizando también para la EP, pudiendo realzar y prolongar la respuesta de la levopoda. También se está experimentando con injertos naturales o el trasplante de células nerviosas, así como con el uso de células genéticamente producidas.\\
		
		También es preciso un abordaje integral y multidisciplinar de la enfermedad, con el objetivo de mejorar o mantener la calidad de vida del paciente y sus cuidadores. Desde el ejercicio físico, la fisioterapia, las terapias ocupacionales, entre otras. Todas ellas suman para retrasar el deterioro.		
		
		\subsubsection{Esclerosis lateral amiotrófica (ELA)}
		
		La ELA es la enfermedad de la neurona motora más frecuente en adultos y se basa en la degeneración combinada de las motoneuronas superiores e inferiores. La consecuencia es una debilidad muscular que avanza hasta la parálisis. Su incidencia está entre 0,4-1,8/100000hab y año. Además, es más frecuente en hombres y la supervivencia media es de 3 años.\\
		
		Los primeros síntomas suelen ser espasmos musculares y debilidad en un miembro o dificultad para hablar, afectando posteriormente a los músculos necesarios para moverse, hablar o comer. Sólo un 10 \% puede desarrollar deterioro cognitivo. Por otra parte, es difícil de diagnosticar de manera temprana y un método de ayuda es el electromiograma. Actualmente no existe ningún tratamiento eficaz, el único que ha demostrado prolongar la supervivencia es el riluzol. Todos los esfuerzos deben ir encaminados a conseguir mejor calidad de vida de pacientes y familiares. El avance de la enfermedad plantea decisiones de gran dificultad ética que deben ser compartidos con los pacientes y su familia.
		
		
		\subsubsection{Enfermedad de Huntington}
		Se trata de un trastorno neurodegenerativo con carácter hereditario autosómico dominante. Presenta tres manifestaciones clínicas principales: movimientos involuntarios bruscos, alteraciones de la personalidad y deterioro mental. Con una prevalencia de 5-10 casos por 100.000 habitantes, afecta igual a hombres y mujeres.\\
		
		El diagnóstico se basa en las características clínicas y los antecedentes familiares confirmados por análisis genético del ADN. No existe tratamiento que cure o ralentice el avance, por lo que está orientado a mejorar la calidad de vida a través del tratamiento sintomático. Las vías de investigación incluyen la interferencia en RNA que operaría en el inicio del proceso. La asistencia sociosanitaria es clave en este trastorno.
				
	\subsection{Intervención terapéutica}

	Los condicionantes del estilo de vida tienen un profundo impacto en nuestro cerebro. Son tres los elementos principales a tener en consideración para prevenir los procesos neurodegenerativos: ejercicio físico, nutrición y actividad intelectual.\\
	
	Hoy podemos establecer que la actividad física es uno de los mejores neuroprotectores existentes, demostrando mejores parámetros cognitivos asociados a la edad. Se sabe que el ejercicio físico influye en la capacidad de generar nuevas neuronas y nuevas conexiones, favoreciendo un incremento de las neurotrofinas que promueven mayores conexiones entre áreas del cerebro, además de su acción antiinflamatoria. Existe diferentes tipos de ejercicio (aeróbico, anaeróbico y de resistencia) que deben tenerse en cuenta en función de estado del paciente, en base a un análisis individualizado. Como conclusión, el ejercicio físico mejora la función cognitiva en adultos mayores de 50 años, debiendo incluir una duración mínima de 45 minutos en intensidad de moderada a vigorosa.\\
	
	Por otro lado, hay que señalar que a actividad intelectual está directamente relacionada con el concepto de reserva cognitiva, como la capacidad del cerebro para tolerar diversas patologías. Existen varios factores asociados al desarrollo de nuestra capacidad cognitiva: volumen cerebral, educación y complejidad laboral, bilingüismo y actividad mental y estimulación cognitiva. Sin embargo, los estudios sobre estos aspectos están sujetos a limitaciones metodológicas, siendo necesario el desarrollo de estudios multidisciplinares que avancen en la evidencia científica.\\
	
	Por último, no podemos dejar de hablar de la importancia del consumo de alimentos saludables como un gran aliado en la prevención de enfermedades neurodegenerativas. Diferentes estudios señalan que modificaciones nutricionales pueden tener un fuerte impacto en la salud. Más allá del papel que puedan desempeñar ciertos nutrientes, lo más relevante es la adopción de la dieta mediterránea como estilo de vida saludable. Los componentes de esta pueden actuar como factores de neuroprotección, como restricción calórica o como antioxidantes naturales. Hay pocos estudios que relacionen patrones dietéticos y enfermedades neurodegenerativas, centrándose más en alimentos aislados. Quedan preguntas por contestar, pero parece claro que la intervención deber realizarse incluso antes de la edad adulta para conseguir la mayor acción preventiva.
	
\end{document}